% --- Archon physics formalization source begin ---
% archon:physics
% archon:covers PhyXMiniProblems/problem_phyx_mini_0605.lean
% archon:source-report reports/phyx_mini/problem_phyx_mini_0605.source.json

\chapter{Physics problem phyx\_mini\_0605}
\label{ch:PhyXMiniProblems_problem_phyx_mini_0605}

\paragraph{Problem source.}
\#\# Physical scenario

Two particles (masses \$m\_1\$ and \$m\_2\$) are attached to the ends of a massless rigid rod of length \$a\$. The system is free to rotate in three dimensions about the (fixed) center of mass.Figure shows a portion of the rotational spectrum of carbon monoxide (CO).

\#\# Question

What is the frequency separation \} (\textbackslash{}Delta 
u) \textbackslash{}text\{ between adjacent lines? Look up the masses of \}\textasciicircum{}\{12\}\textbackslash{}text\{C and \}\textasciicircum{}\{16\}\textbackslash{}text\{O, and from \} m\_1, m\_2, \textbackslash{}text\{ and \} \textbackslash{}Delta 
u \textbackslash{}text\{ determine the distance between the atoms.\}

\#\# Answer choices

- A: A: 1.6 \textbackslash{}times 10\textasciicircum{}\{-10\} \textbackslash{}

- B: B: 2.2 \textbackslash{}times 10\textasciicircum{}\{-10\} \textbackslash{}

- C: C: 1.1 \textbackslash{}times 10\textasciicircum{}\{-10\} \textbackslash{}

- D: D: 7.5 \textbackslash{}times 10\textasciicircum{}\{-11\} \textbackslash{}

\#\# Figure caption (auxiliary; use the image as primary evidence)

The image is a graph displaying a spectral pattern or a waveform. Here are the details:

- **Horizontal Axis (X-Axis):** Labeled with numerical values at intervals; visible values are 10, 30, 50, and 70.
  
- **Vertical Axis (Y-Axis):** Not explicitly labeled, but it measures the amplitude or intensity of the waveform.

- **Waveform:** The graph shows a complex line with peaks and troughs, indicating varying intensity values. It starts with a gradual increase, reaches a series of high peaks around the middle, and then gradually decreases towards the end.

- **Overall Shape:** The waveform rises from the left, reaches a high point with multiple prominent peaks, and then decreases towards the right.

There are no titles, labels, or additional text visible on the axes, but the data suggests analysis typical for spectral or frequency-related studies.

\#\# Dataset metadata

Quantum Phenomena; Physical Model Grounding Reasoning, Multi-Formula Reasoning

\paragraph{Recorded answer/context.}
C: C: 1.1 \textbackslash{}times 10\textasciicircum{}\{-10\} \textbackslash{}

\paragraph{Figure/image path.}
/root/proposal\_for\_physic/science-mango/phyx\_mini\_run/phyx\_data/test\_image/605.png

\paragraph{Formalization target.}
create a compiling Lean file with sorry bodies at `PhyXMiniProblems/problem\_phyx\_mini\_0605.lean`. The Lean declarations must preserve the physical quantities, dimensions or dimensional roles, figure labels, governing-law hypotheses, and final relation expressed by this problem.
Use Mathlib/Physlib names found through LeanExplore where available. If a physics API is missing, introduce faithful local abstractions rather than scalar placeholder aliases.

\begin{theorem}[Physics formalization target]
\label{thm:physics:phyx_mini_0605:target}
\lean{PhyXMiniProblems.ProblemPhyXMini0605.carbonMonoxideSpectrum_determines_frequencySeparation_and_choiceC}
\uses{def:physics:phyx-mini-0605:phyxminiproblems-problemphyxmini0605-carbonmonoxiderotorsetup, def:physics:phyx-mini-0605:phyxminiproblems-problemphyxmini0605-matchescarbonmonoxiderotorscenario, def:physics:phyx-mini-0605:phyxminiproblems-problemphyxmini0605-matchesprimarycarbonmonoxidespectrum, def:physics:phyx-mini-0605:phyxminiproblems-problemphyxmini0605-usescarbonmonoxidereferencedata, def:physics:phyx-mini-0605:phyxminiproblems-problemphyxmini0605-hasphysicalcarbonmonoxideparameters, def:physics:phyx-mini-0605:phyxminiproblems-problemphyxmini0605-satisfiestwobodycenterofmassgeometry, def:physics:phyx-mini-0605:phyxminiproblems-problemphyxmini0605-satisfiescarbonmonoxiderigidrotorlaws, def:physics:phyx-mini-0605:phyxminiproblems-problemphyxmini0605-agreeswithfigurefrequencyseparation, def:physics:phyx-mini-0605:phyxminiproblems-problemphyxmini0605-agreeswithdisplayedbondlength, def:physics:phyx-mini-0605:phyxminiproblems-problemphyxmini0605-isuniqueclosestbondlengthchoice}
The assigned autoformalize agent should translate this physics problem into Lean declarations in the covered file, with theorem and lemma proof bodies written as `by sorry`.
\end{theorem}
\begin{proof}
This is an autoformalization task, not a proof task. Produce statements that can later be proved by the physics prover without weakening the physical model.
\end{proof}

% --- Archon declaration topology begin ---
\section{Declaration topology}
\begin{definition}[Length Quantity]
  \label{def:physics:phyx-mini-0605:phyxminiproblems-problemphyxmini0605-lengthquantity}
  \lean{PhyXMiniProblems.ProblemPhyXMini0605.LengthQuantity}
  A nonnegative, unit-independent physical length
\end{definition}
\begin{proof}
  This is the declared model definition of Length Quantity.
\end{proof}
\begin{definition}[Mass Quantity]
  \label{def:physics:phyx-mini-0605:phyxminiproblems-problemphyxmini0605-massquantity}
  \lean{PhyXMiniProblems.ProblemPhyXMini0605.MassQuantity}
  A nonnegative, unit-independent physical mass
\end{definition}
\begin{proof}
  This is the declared model definition of Mass Quantity.
\end{proof}
\begin{definition}[moment Of Inertia Dimension]
  \label{def:physics:phyx-mini-0605:phyxminiproblems-problemphyxmini0605-momentofinertiadimension}
  \lean{PhyXMiniProblems.ProblemPhyXMini0605.momentOfInertiaDimension}
  The physical dimension of a moment of inertia, mass * length²
\end{definition}
\begin{proof}
  This is the declared model definition of moment Of Inertia Dimension.
\end{proof}
\begin{definition}[Moment Of Inertia Quantity]
  \label{def:physics:phyx-mini-0605:phyxminiproblems-problemphyxmini0605-momentofinertiaquantity}
  \lean{PhyXMiniProblems.ProblemPhyXMini0605.MomentOfInertiaQuantity}
  \uses{def:physics:phyx-mini-0605:phyxminiproblems-problemphyxmini0605-momentofinertiadimension}
  A nonnegative, unit-independent moment of inertia
\end{definition}
\begin{proof}
  This is the declared model definition of Moment Of Inertia Quantity.
\end{proof}
\begin{definition}[action Dimension]
  \label{def:physics:phyx-mini-0605:phyxminiproblems-problemphyxmini0605-actiondimension}
  \lean{PhyXMiniProblems.ProblemPhyXMini0605.actionDimension}
  The physical dimension of action, equivalently energy times time
\end{definition}
\begin{proof}
  This is the declared model definition of action Dimension.
\end{proof}
\begin{definition}[Action Quantity]
  \label{def:physics:phyx-mini-0605:phyxminiproblems-problemphyxmini0605-actionquantity}
  \lean{PhyXMiniProblems.ProblemPhyXMini0605.ActionQuantity}
  \uses{def:physics:phyx-mini-0605:phyxminiproblems-problemphyxmini0605-actiondimension}
  A nonnegative, unit-independent physical action
\end{definition}
\begin{proof}
  This is the declared model definition of Action Quantity.
\end{proof}
\begin{definition}[Frequency Quantity]
  \label{def:physics:phyx-mini-0605:phyxminiproblems-problemphyxmini0605-frequencyquantity}
  \lean{PhyXMiniProblems.ProblemPhyXMini0605.FrequencyQuantity}
  A nonnegative ordinary frequency, with dimension inverse time
\end{definition}
\begin{proof}
  This is the declared model definition of Frequency Quantity.
\end{proof}
\begin{definition}[Wavenumber Quantity]
  \label{def:physics:phyx-mini-0605:phyxminiproblems-problemphyxmini0605-wavenumberquantity}
  \lean{PhyXMiniProblems.ProblemPhyXMini0605.WavenumberQuantity}
  A nonnegative spectroscopic wavenumber, with dimension inverse length
\end{definition}
\begin{proof}
  This is the declared model definition of Wavenumber Quantity.
\end{proof}
\begin{definition}[Speed Quantity]
  \label{def:physics:phyx-mini-0605:phyxminiproblems-problemphyxmini0605-speedquantity}
  \lean{PhyXMiniProblems.ProblemPhyXMini0605.SpeedQuantity}
  A unit-independent speed with a real carrier. This is the exact type of Physlib's dimensionful constant DimSpeed.speedOfLight
\end{definition}
\begin{proof}
  This is the declared model definition of Speed Quantity.
\end{proof}
\begin{definition}[length In Meters]
  \label{def:physics:phyx-mini-0605:phyxminiproblems-problemphyxmini0605-lengthinmeters}
  \lean{PhyXMiniProblems.ProblemPhyXMini0605.lengthInMeters}
  \uses{def:physics:phyx-mini-0605:phyxminiproblems-problemphyxmini0605-lengthquantity}
  Metre readout of a physical length
\end{definition}
\begin{proof}
  This is the declared model definition of length In Meters.
\end{proof}
\begin{definition}[mass In Kilograms]
  \label{def:physics:phyx-mini-0605:phyxminiproblems-problemphyxmini0605-massinkilograms}
  \lean{PhyXMiniProblems.ProblemPhyXMini0605.massInKilograms}
  \uses{def:physics:phyx-mini-0605:phyxminiproblems-problemphyxmini0605-massquantity}
  Kilogram readout of a physical mass
\end{definition}
\begin{proof}
  This is the declared model definition of mass In Kilograms.
\end{proof}
\begin{definition}[mass In Atomic Mass Units]
  \label{def:physics:phyx-mini-0605:phyxminiproblems-problemphyxmini0605-massinatomicmassunits}
  \lean{PhyXMiniProblems.ProblemPhyXMini0605.massInAtomicMassUnits}
  \uses{def:physics:phyx-mini-0605:phyxminiproblems-problemphyxmini0605-massquantity, def:physics:phyx-mini-0605:phyxminiproblems-problemphyxmini0605-massinkilograms}
  A mass expressed as a dimensionless multiple of one atomic mass unit
\end{definition}
\begin{proof}
  This is the declared model definition of mass In Atomic Mass Units.
\end{proof}
\begin{definition}[moment Of Inertia In Kilogram Meters Squared]
  \label{def:physics:phyx-mini-0605:phyxminiproblems-problemphyxmini0605-momentofinertiainkilogrammeterssquared}
  \lean{PhyXMiniProblems.ProblemPhyXMini0605.momentOfInertiaInKilogramMetersSquared}
  \uses{def:physics:phyx-mini-0605:phyxminiproblems-problemphyxmini0605-momentofinertiaquantity}
  Coherent-SI readout of a moment of inertia, in kg m²
\end{definition}
\begin{proof}
  This is the declared model definition of moment Of Inertia In Kilogram Meters Squared.
\end{proof}
\begin{definition}[action In Joule Seconds]
  \label{def:physics:phyx-mini-0605:phyxminiproblems-problemphyxmini0605-actioninjouleseconds}
  \lean{PhyXMiniProblems.ProblemPhyXMini0605.actionInJouleSeconds}
  \uses{def:physics:phyx-mini-0605:phyxminiproblems-problemphyxmini0605-actionquantity}
  Coherent-SI readout of an action, in joule-seconds
\end{definition}
\begin{proof}
  This is the declared model definition of action In Joule Seconds.
\end{proof}
\begin{definition}[energy In Joules]
  \label{def:physics:phyx-mini-0605:phyxminiproblems-problemphyxmini0605-energyinjoules}
  \lean{PhyXMiniProblems.ProblemPhyXMini0605.energyInJoules}
  Coherent-SI readout of an energy, in joules
\end{definition}
\begin{proof}
  This is the declared model definition of energy In Joules.
\end{proof}
\begin{definition}[frequency In Hertz]
  \label{def:physics:phyx-mini-0605:phyxminiproblems-problemphyxmini0605-frequencyinhertz}
  \lean{PhyXMiniProblems.ProblemPhyXMini0605.frequencyInHertz}
  \uses{def:physics:phyx-mini-0605:phyxminiproblems-problemphyxmini0605-frequencyquantity}
  Coherent-SI readout of an ordinary frequency, in hertz
\end{definition}
\begin{proof}
  This is the declared model definition of frequency In Hertz.
\end{proof}
\begin{definition}[wavenumber In Inverse Centimeters]
  \label{def:physics:phyx-mini-0605:phyxminiproblems-problemphyxmini0605-wavenumberininversecentimeters}
  \lean{PhyXMiniProblems.ProblemPhyXMini0605.wavenumberInInverseCentimeters}
  \uses{def:physics:phyx-mini-0605:phyxminiproblems-problemphyxmini0605-wavenumberquantity}
  Wavenumber readout in inverse centimeters
\end{definition}
\begin{proof}
  This is the declared model definition of wavenumber In Inverse Centimeters.
\end{proof}
\begin{definition}[speed In Meters Per Second]
  \label{def:physics:phyx-mini-0605:phyxminiproblems-problemphyxmini0605-speedinmeterspersecond}
  \lean{PhyXMiniProblems.ProblemPhyXMini0605.speedInMetersPerSecond}
  \uses{def:physics:phyx-mini-0605:phyxminiproblems-problemphyxmini0605-speedquantity}
  Coherent-SI speed readout, in metres per second
\end{definition}
\begin{proof}
  This is the declared model definition of speed In Meters Per Second.
\end{proof}
\begin{definition}[Atom Site]
  \label{def:physics:phyx-mini-0605:phyxminiproblems-problemphyxmini0605-atomsite}
  \lean{PhyXMiniProblems.ProblemPhyXMini0605.AtomSite}
  The two atomic sites at the ends of the massless rod
\end{definition}
\begin{proof}
  This is the declared model definition of Atom Site.
\end{proof}
\begin{definition}[Isotope Label]
  \label{def:physics:phyx-mini-0605:phyxminiproblems-problemphyxmini0605-isotopelabel}
  \lean{PhyXMiniProblems.ProblemPhyXMini0605.IsotopeLabel}
  Isotopes specified by the problem
\end{definition}
\begin{proof}
  This is the declared model definition of Isotope Label.
\end{proof}
\begin{definition}[Bond Model]
  \label{def:physics:phyx-mini-0605:phyxminiproblems-problemphyxmini0605-bondmodel}
  \lean{PhyXMiniProblems.ProblemPhyXMini0605.BondModel}
  Mechanical model assigned to the bond joining the atoms
\end{definition}
\begin{proof}
  This is the declared model definition of Bond Model.
\end{proof}
\begin{definition}[Rotation Center]
  \label{def:physics:phyx-mini-0605:phyxminiproblems-problemphyxmini0605-rotationcenter}
  \lean{PhyXMiniProblems.ProblemPhyXMini0605.RotationCenter}
  Point about which the molecular system rotates
\end{definition}
\begin{proof}
  This is the declared model definition of Rotation Center.
\end{proof}
\begin{definition}[Rotation Freedom]
  \label{def:physics:phyx-mini-0605:phyxminiproblems-problemphyxmini0605-rotationfreedom}
  \lean{PhyXMiniProblems.ProblemPhyXMini0605.RotationFreedom}
  Spatial freedom of the molecular rotation
\end{definition}
\begin{proof}
  This is the declared model definition of Rotation Freedom.
\end{proof}
\begin{definition}[Figure Axis]
  \label{def:physics:phyx-mini-0605:phyxminiproblems-problemphyxmini0605-figureaxis}
  \lean{PhyXMiniProblems.ProblemPhyXMini0605.FigureAxis}
  The two axes visible in image 605
\end{definition}
\begin{proof}
  This is the declared model definition of Figure Axis.
\end{proof}
\begin{definition}[Spectrum Axis Quantity]
  \label{def:physics:phyx-mini-0605:phyxminiproblems-problemphyxmini0605-spectrumaxisquantity}
  \lean{PhyXMiniProblems.ProblemPhyXMini0605.SpectrumAxisQuantity}
  Physical interpretation of each plotted axis
\end{definition}
\begin{proof}
  This is the declared model definition of Spectrum Axis Quantity.
\end{proof}
\begin{definition}[Spectrum Axis Unit]
  \label{def:physics:phyx-mini-0605:phyxminiproblems-problemphyxmini0605-spectrumaxisunit}
  \lean{PhyXMiniProblems.ProblemPhyXMini0605.SpectrumAxisUnit}
  Unit convention assigned to each plotted axis
\end{definition}
\begin{proof}
  This is the declared model definition of Spectrum Axis Unit.
\end{proof}
\begin{definition}[Horizontal Tick]
  \label{def:physics:phyx-mini-0605:phyxminiproblems-problemphyxmini0605-horizontaltick}
  \lean{PhyXMiniProblems.ProblemPhyXMini0605.HorizontalTick}
  The four numerical horizontal-axis ticks visible in the raster
\end{definition}
\begin{proof}
  This is the declared model definition of Horizontal Tick.
\end{proof}
\begin{definition}[Representative Spectrum Line]
  \label{def:physics:phyx-mini-0605:phyxminiproblems-problemphyxmini0605-representativespectrumline}
  \lean{PhyXMiniProblems.ProblemPhyXMini0605.RepresentativeSpectrumLine}
  Three consecutive central absorption troughs used for the spacing readout
\end{definition}
\begin{proof}
  This is the declared model definition of Representative Spectrum Line.
\end{proof}
\begin{definition}[Carbon Monoxide Spectrum Figure]
  \label{def:physics:phyx-mini-0605:phyxminiproblems-problemphyxmini0605-carbonmonoxidespectrumfigure}
  \lean{PhyXMiniProblems.ProblemPhyXMini0605.CarbonMonoxideSpectrumFigure}
  \uses{def:physics:phyx-mini-0605:phyxminiproblems-problemphyxmini0605-wavenumberquantity, def:physics:phyx-mini-0605:phyxminiproblems-problemphyxmini0605-figureaxis, def:physics:phyx-mini-0605:phyxminiproblems-problemphyxmini0605-spectrumaxisquantity, def:physics:phyx-mini-0605:phyxminiproblems-problemphyxmini0605-spectrumaxisunit, def:physics:phyx-mini-0605:phyxminiproblems-problemphyxmini0605-horizontaltick, def:physics:phyx-mini-0605:phyxminiproblems-problemphyxmini0605-representativespectrumline}
  Data carried by the supplied spectrum image. The wavenumbers are physical quantities; the remaining scalar fields are literal raster labels or display features. Frequencies are stored separately in the physical setup below
\end{definition}
\begin{proof}
  This is the declared model definition of Carbon Monoxide Spectrum Figure.
\end{proof}
\begin{definition}[Carbon Monoxide Rotor Setup]
  \label{def:physics:phyx-mini-0605:phyxminiproblems-problemphyxmini0605-carbonmonoxiderotorsetup}
  \lean{PhyXMiniProblems.ProblemPhyXMini0605.CarbonMonoxideRotorSetup}
  \uses{def:physics:phyx-mini-0605:phyxminiproblems-problemphyxmini0605-lengthquantity, def:physics:phyx-mini-0605:phyxminiproblems-problemphyxmini0605-massquantity, def:physics:phyx-mini-0605:phyxminiproblems-problemphyxmini0605-momentofinertiaquantity, def:physics:phyx-mini-0605:phyxminiproblems-problemphyxmini0605-actionquantity, def:physics:phyx-mini-0605:phyxminiproblems-problemphyxmini0605-frequencyquantity, def:physics:phyx-mini-0605:phyxminiproblems-problemphyxmini0605-speedquantity, def:physics:phyx-mini-0605:phyxminiproblems-problemphyxmini0605-atomsite, def:physics:phyx-mini-0605:phyxminiproblems-problemphyxmini0605-isotopelabel, def:physics:phyx-mini-0605:phyxminiproblems-problemphyxmini0605-bondmodel, def:physics:phyx-mini-0605:phyxminiproblems-problemphyxmini0605-rotationcenter, def:physics:phyx-mini-0605:phyxminiproblems-problemphyxmini0605-rotationfreedom, def:physics:phyx-mini-0605:phyxminiproblems-problemphyxmini0605-representativespectrumline, def:physics:phyx-mini-0605:phyxminiproblems-problemphyxmini0605-carbonmonoxidespectrumfigure}
  The physical carbon-monoxide rotor and the observables appearing in its spectrum. No field is defined from a displayed answer. The line frequencies, their common separation, the moment of inertia, and the bond length are related only by the governing-law structures below
\end{definition}
\begin{proof}
  This is the declared model definition of Carbon Monoxide Rotor Setup.
\end{proof}
\begin{definition}[Matches Carbon Monoxide Rotor Scenario]
  \label{def:physics:phyx-mini-0605:phyxminiproblems-problemphyxmini0605-matchescarbonmonoxiderotorscenario}
  \lean{PhyXMiniProblems.ProblemPhyXMini0605.MatchesCarbonMonoxideRotorScenario}
  \uses{def:physics:phyx-mini-0605:phyxminiproblems-problemphyxmini0605-carbonmonoxiderotorsetup}
  Qualitative model stated in the problem prose
\end{definition}
\begin{proof}
  This is the declared model definition of Matches Carbon Monoxide Rotor Scenario.
\end{proof}
\begin{definition}[Matches Primary Carbon Monoxide Spectrum]
  \label{def:physics:phyx-mini-0605:phyxminiproblems-problemphyxmini0605-matchesprimarycarbonmonoxidespectrum}
  \lean{PhyXMiniProblems.ProblemPhyXMini0605.MatchesPrimaryCarbonMonoxideSpectrum}
  \uses{def:physics:phyx-mini-0605:phyxminiproblems-problemphyxmini0605-wavenumberininversecentimeters, def:physics:phyx-mini-0605:phyxminiproblems-problemphyxmini0605-carbonmonoxiderotorsetup}
  Literal and calibrated readouts from image 605. The ticks at 10, 30, 50, and 70 are visible. Three central consecutive troughs lie, to the precision supported by the raster, at approximately 42, 46, and 50 cm⁻¹. These are wavenumber data, not a premise about frequency separation or bond length
\end{definition}
\begin{proof}
  This is the declared model definition of Matches Primary Carbon Monoxide Spectrum.
\end{proof}
\begin{definition}[Uses Carbon Monoxide Reference Data]
  \label{def:physics:phyx-mini-0605:phyxminiproblems-problemphyxmini0605-usescarbonmonoxidereferencedata}
  \lean{PhyXMiniProblems.ProblemPhyXMini0605.UsesCarbonMonoxideReferenceData}
  \uses{def:physics:phyx-mini-0605:phyxminiproblems-problemphyxmini0605-massinkilograms, def:physics:phyx-mini-0605:phyxminiproblems-problemphyxmini0605-massinatomicmassunits, def:physics:phyx-mini-0605:phyxminiproblems-problemphyxmini0605-actioninjouleseconds, def:physics:phyx-mini-0605:phyxminiproblems-problemphyxmini0605-speedinmeterspersecond, def:physics:phyx-mini-0605:phyxminiproblems-problemphyxmini0605-carbonmonoxiderotorsetup}
  Looked-up isotope data and calibrated universal constants. Physlib grounds the reduced Planck constant Constants.ℏ and the dimensionful speed of light DimSpeed.speedOfLight. No target separation or bond length occurs here
\end{definition}
\begin{proof}
  This is the declared model definition of Uses Carbon Monoxide Reference Data.
\end{proof}
\begin{definition}[Has Physical Carbon Monoxide Parameters]
  \label{def:physics:phyx-mini-0605:phyxminiproblems-problemphyxmini0605-hasphysicalcarbonmonoxideparameters}
  \lean{PhyXMiniProblems.ProblemPhyXMini0605.HasPhysicalCarbonMonoxideParameters}
  \uses{def:physics:phyx-mini-0605:phyxminiproblems-problemphyxmini0605-lengthinmeters, def:physics:phyx-mini-0605:phyxminiproblems-problemphyxmini0605-massinkilograms, def:physics:phyx-mini-0605:phyxminiproblems-problemphyxmini0605-momentofinertiainkilogrammeterssquared, def:physics:phyx-mini-0605:phyxminiproblems-problemphyxmini0605-actioninjouleseconds, def:physics:phyx-mini-0605:phyxminiproblems-problemphyxmini0605-energyinjoules, def:physics:phyx-mini-0605:phyxminiproblems-problemphyxmini0605-frequencyinhertz, def:physics:phyx-mini-0605:phyxminiproblems-problemphyxmini0605-speedinmeterspersecond, def:physics:phyx-mini-0605:phyxminiproblems-problemphyxmini0605-carbonmonoxiderotorsetup}
  Positivity and nondegeneracy conditions selecting the physical branch
\end{definition}
\begin{proof}
  This is the declared model definition of Has Physical Carbon Monoxide Parameters.
\end{proof}
\begin{definition}[Satisfies Two Body Center Of Mass Geometry]
  \label{def:physics:phyx-mini-0605:phyxminiproblems-problemphyxmini0605-satisfiestwobodycenterofmassgeometry}
  \lean{PhyXMiniProblems.ProblemPhyXMini0605.SatisfiesTwoBodyCenterOfMassGeometry}
  \uses{def:physics:phyx-mini-0605:phyxminiproblems-problemphyxmini0605-lengthinmeters, def:physics:phyx-mini-0605:phyxminiproblems-problemphyxmini0605-massinkilograms, def:physics:phyx-mini-0605:phyxminiproblems-problemphyxmini0605-carbonmonoxiderotorsetup}
  Center-of-mass geometry for two point masses on opposite ends of the rod. The two lever arms add to the interatomic distance and their mass moments balance
\end{definition}
\begin{proof}
  This is the declared model definition of Satisfies Two Body Center Of Mass Geometry.
\end{proof}
\begin{definition}[Satisfies Carbon Monoxide Rigid Rotor Laws]
  \label{def:physics:phyx-mini-0605:phyxminiproblems-problemphyxmini0605-satisfiescarbonmonoxiderigidrotorlaws}
  \lean{PhyXMiniProblems.ProblemPhyXMini0605.SatisfiesCarbonMonoxideRigidRotorLaws}
  \uses{def:physics:phyx-mini-0605:phyxminiproblems-problemphyxmini0605-lengthinmeters, def:physics:phyx-mini-0605:phyxminiproblems-problemphyxmini0605-massinkilograms, def:physics:phyx-mini-0605:phyxminiproblems-problemphyxmini0605-momentofinertiainkilogrammeterssquared, def:physics:phyx-mini-0605:phyxminiproblems-problemphyxmini0605-actioninjouleseconds, def:physics:phyx-mini-0605:phyxminiproblems-problemphyxmini0605-energyinjoules, def:physics:phyx-mini-0605:phyxminiproblems-problemphyxmini0605-frequencyinhertz, def:physics:phyx-mini-0605:phyxminiproblems-problemphyxmini0605-wavenumberininversecentimeters, def:physics:phyx-mini-0605:phyxminiproblems-problemphyxmini0605-speedinmeterspersecond, def:physics:phyx-mini-0605:phyxminiproblems-problemphyxmini0605-carbonmonoxiderotorsetup}
  Rigid-rotor and spectrometer laws: I = m\_C r\_C² + m\_O r\_O²; μ (m\_C + m\_O) = m\_C m\_O; E\_l = ℏ² l(l+1)/(2I); h ν\_l = E\_\{l+1\} - E\_l, with h = 2πℏ; the three representative troughs are consecutive rotor transitions; ν = c times wavenumber, including 100 m⁻¹ per cm⁻¹; and the independent separation observable is the difference of adjacent lines. These are general physical and calibration relations. None inserts a numeric frequency separation or interatomic distance
\end{definition}
\begin{proof}
  This is the declared model definition of Satisfies Carbon Monoxide Rigid Rotor Laws.
\end{proof}
\begin{lemma}[representative Wavenumber Spacing eq four]
  \label{lem:physics:phyx-mini-0605:phyxminiproblems-problemphyxmini0605-representativewavenumberspacing-eq-four}
  \lean{PhyXMiniProblems.ProblemPhyXMini0605.representativeWavenumberSpacing_eq_four}
  \uses{def:physics:phyx-mini-0605:phyxminiproblems-problemphyxmini0605-wavenumberininversecentimeters, def:physics:phyx-mini-0605:phyxminiproblems-problemphyxmini0605-carbonmonoxiderotorsetup, def:physics:phyx-mini-0605:phyxminiproblems-problemphyxmini0605-matchesprimarycarbonmonoxidespectrum}
  The calibrated central troughs are separated by 4 cm⁻¹
\end{lemma}
\begin{proof}
  The conclusion follows from the governing relations and figure readouts named in the statement of representative Wavenumber Spacing eq four.
\end{proof}
\begin{lemma}[moment Of Inertia eq reduced Mass mul bond Length sq]
  \label{lem:physics:phyx-mini-0605:phyxminiproblems-problemphyxmini0605-momentofinertia-eq-reducedmass-mul-bondlength-sq}
  \lean{PhyXMiniProblems.ProblemPhyXMini0605.momentOfInertia_eq_reducedMass_mul_bondLength_sq}
  \uses{def:physics:phyx-mini-0605:phyxminiproblems-problemphyxmini0605-lengthinmeters, def:physics:phyx-mini-0605:phyxminiproblems-problemphyxmini0605-massinkilograms, def:physics:phyx-mini-0605:phyxminiproblems-problemphyxmini0605-momentofinertiainkilogrammeterssquared, def:physics:phyx-mini-0605:phyxminiproblems-problemphyxmini0605-carbonmonoxiderotorsetup, def:physics:phyx-mini-0605:phyxminiproblems-problemphyxmini0605-hasphysicalcarbonmonoxideparameters, def:physics:phyx-mini-0605:phyxminiproblems-problemphyxmini0605-satisfiestwobodycenterofmassgeometry, def:physics:phyx-mini-0605:phyxminiproblems-problemphyxmini0605-satisfiescarbonmonoxiderigidrotorlaws}
  For two point masses at their common center of mass, the point-mass inertia reduces to I = μ a²
\end{lemma}
\begin{proof}
  The conclusion follows from the governing relations and figure readouts named in the statement of moment Of Inertia eq reduced Mass mul bond Length sq.
\end{proof}
\begin{definition}[Agrees With Figure Frequency Separation]
  \label{def:physics:phyx-mini-0605:phyxminiproblems-problemphyxmini0605-agreeswithfigurefrequencyseparation}
  \lean{PhyXMiniProblems.ProblemPhyXMini0605.AgreesWithFigureFrequencySeparation}
  \uses{def:physics:phyx-mini-0605:phyxminiproblems-problemphyxmini0605-frequencyinhertz, def:physics:phyx-mini-0605:phyxminiproblems-problemphyxmini0605-carbonmonoxiderotorsetup}
  The figure-derived adjacent-line separation agrees with 1.20 × 10¹¹ Hz to within 10⁹ Hz
\end{definition}
\begin{proof}
  This is the declared model definition of Agrees With Figure Frequency Separation.
\end{proof}
\begin{lemma}[adjacent Frequency Separation agrees with figure]
  \label{lem:physics:phyx-mini-0605:phyxminiproblems-problemphyxmini0605-adjacentfrequencyseparation-agrees-with-figure}
  \lean{PhyXMiniProblems.ProblemPhyXMini0605.adjacentFrequencySeparation_agrees_with_figure}
  \uses{def:physics:phyx-mini-0605:phyxminiproblems-problemphyxmini0605-carbonmonoxiderotorsetup, def:physics:phyx-mini-0605:phyxminiproblems-problemphyxmini0605-matchesprimarycarbonmonoxidespectrum, def:physics:phyx-mini-0605:phyxminiproblems-problemphyxmini0605-usescarbonmonoxidereferencedata, def:physics:phyx-mini-0605:phyxminiproblems-problemphyxmini0605-satisfiescarbonmonoxiderigidrotorlaws, def:physics:phyx-mini-0605:phyxminiproblems-problemphyxmini0605-agreeswithfigurefrequencyseparation}
  The calibrated spectrum determines the stated hertz-scale separation
\end{lemma}
\begin{proof}
  The conclusion follows from the governing relations and figure readouts named in the statement of adjacent Frequency Separation agrees with figure.
\end{proof}
\begin{definition}[Answer Choice]
  \label{def:physics:phyx-mini-0605:phyxminiproblems-problemphyxmini0605-answerchoice}
  \lean{PhyXMiniProblems.ProblemPhyXMini0605.AnswerChoice}
  Labels of the four interatomic-distance choices
\end{definition}
\begin{proof}
  This is the declared model definition of Answer Choice.
\end{proof}
\begin{definition}[bond Length Meters]
  \label{def:physics:phyx-mini-0605:phyxminiproblems-problemphyxmini0605-answerchoice-bondlengthmeters}
  \lean{PhyXMiniProblems.ProblemPhyXMini0605.AnswerChoice.bondLengthMeters}
  \uses{def:physics:phyx-mini-0605:phyxminiproblems-problemphyxmini0605-answerchoice}
  Bond length printed beside an answer choice, in metres
\end{definition}
\begin{proof}
  This is the declared model definition of bond Length Meters.
\end{proof}
\begin{definition}[recorded Dataset Answer]
  \label{def:physics:phyx-mini-0605:phyxminiproblems-problemphyxmini0605-recordeddatasetanswer}
  \lean{PhyXMiniProblems.ProblemPhyXMini0605.recordedDatasetAnswer}
  \uses{def:physics:phyx-mini-0605:phyxminiproblems-problemphyxmini0605-answerchoice}
  Dataset metadata records choice C; this is not a theorem premise
\end{definition}
\begin{proof}
  This is the declared model definition of recorded Dataset Answer.
\end{proof}
\begin{definition}[answer Choice Error Meters]
  \label{def:physics:phyx-mini-0605:phyxminiproblems-problemphyxmini0605-answerchoiceerrormeters}
  \lean{PhyXMiniProblems.ProblemPhyXMini0605.answerChoiceErrorMeters}
  \uses{def:physics:phyx-mini-0605:phyxminiproblems-problemphyxmini0605-lengthinmeters, def:physics:phyx-mini-0605:phyxminiproblems-problemphyxmini0605-carbonmonoxiderotorsetup, def:physics:phyx-mini-0605:phyxminiproblems-problemphyxmini0605-answerchoice}
  Absolute error between the physical bond length and a displayed choice
\end{definition}
\begin{proof}
  This is the declared model definition of answer Choice Error Meters.
\end{proof}
\begin{definition}[Agrees With Displayed Bond Length]
  \label{def:physics:phyx-mini-0605:phyxminiproblems-problemphyxmini0605-agreeswithdisplayedbondlength}
  \lean{PhyXMiniProblems.ProblemPhyXMini0605.AgreesWithDisplayedBondLength}
  \uses{def:physics:phyx-mini-0605:phyxminiproblems-problemphyxmini0605-carbonmonoxiderotorsetup, def:physics:phyx-mini-0605:phyxminiproblems-problemphyxmini0605-answerchoice, def:physics:phyx-mini-0605:phyxminiproblems-problemphyxmini0605-answerchoiceerrormeters}
  Agreement at the precision of the displayed two-significant-figure bond length, allowing 5 × 10⁻¹² m for graph and isotope-mass readout precision
\end{definition}
\begin{proof}
  This is the declared model definition of Agrees With Displayed Bond Length.
\end{proof}
\begin{definition}[Is Unique Closest Bond Length Choice]
  \label{def:physics:phyx-mini-0605:phyxminiproblems-problemphyxmini0605-isuniqueclosestbondlengthchoice}
  \lean{PhyXMiniProblems.ProblemPhyXMini0605.IsUniqueClosestBondLengthChoice}
  \uses{def:physics:phyx-mini-0605:phyxminiproblems-problemphyxmini0605-carbonmonoxiderotorsetup, def:physics:phyx-mini-0605:phyxminiproblems-problemphyxmini0605-answerchoice, def:physics:phyx-mini-0605:phyxminiproblems-problemphyxmini0605-answerchoiceerrormeters}
  A choice is strictly closer than every other displayed bond length
\end{definition}
\begin{proof}
  This is the declared model definition of Is Unique Closest Bond Length Choice.
\end{proof}
% --- Archon declaration topology end ---
% --- Archon physics formalization source end ---
